%%%%%
%
\section{�chantillonnage et recherche du maximum d'une variable sur une nouvelle base
de temps donn�e par un time\_{bnds}}

\itemstyle{Op\'erateur~:} \operfmt{MaxEch\_tps\_bnds}
\index{MaxEch\_tps\_bnds}

%%%%%OBJET%%%%
\itemstyle{Objet~:}
Cette fonction calcule, pour chaque instant $t$ d'une
base de temps $time_{out}$, 
une variable $X_{out}$, contenant le maximum des
valeurs de  $X_{in}$ entre les bornes donn�es par le $time_{bnds}$ associ� � $time_{out}$.

Aucune supposition n'est faite sur la r�gularit� de
l'�chantillonnage des 2 variables. Il est cependant n�cessaire
que $time_{in}$ (le temps de $X_{in}$) soit "� peu pr�s" tri� par ordre
croissant.


%%%%%ARGUMENTS%%%%
\bigskip
%% Dans ce fichier, les �ventuels accents sont cod�s
%% en ISO-8859-15 (aka latin9)
%% La ligne suivante donne trois exemples 
%% e accent aigu : �, e accent grave : �, e-dans-l'o : �
%%%%%ENTREES%%%%
\begin{tabarg}
\multicolumn{3}{l}{\titretab {Entr\'ees}}\\
\hline
$X_{in}$ & Vecteur, \texttt{float} & Variable d'origine\\
    $toler_{mnq}$ & Coeff., \texttt{Int\_32} & Seuil de calcul du maximum (nb. �chant). A chaque instant de la base de temps d'arriv�e, 
    le maximum est calcul� si le nombre d'�chantillons dans l'intervalle de temps est strictement sup�rieur � $toler_{mnq}$. Sinon, 
le r�sultat est "valeur manquante" Si $toler_{mnq}$ = 0, le calcul est fait d�s qu'il y a au moins une valeur valide dans l'�chantillon. \\
$time_{out}$ & Vecteur, \texttt{double} & Base de temps d'arriv�e (secondes)\\
\end{tabarg}

\medskip
%%%%%SORTIES%%%%
\begin{tabarg}
\multicolumn{3}{l}{\titretab {Sorties}}\\
\hline
$X_{out}$ & Vecteur, \texttt{float} & $X_{in}$ maximum �chantillonn� sur
la base de $time_{out}$ (m�me unit� que $X_{in})$\\
\end{tabarg}



%%%%%FORMULE%%%%
\itemstyle{Formule}

%%Max
Recherche de la valeur maximale dans chaque �chantillon
\[
X_M[i] = \max_{{tb[i][0]} \le j< {tb[i][1]}} (X_{in}[j]) 
\]
 
Pour chaque $X_{in}[j]$ diff\'{e}rent de NaN.


%%%%COMMENTAIRE%%%%
\itemstyle{Commentaire~:}
\textit{tb} est un tableau d'intervalles de temps d�termin�s � partir du vecteur temps $time_{out}$.

%%%%%AUTEUR%%%%
\itemstyle{Auteur~:}
CNRM/GMEI/TRAMM

%%%%%REFERENCE%%%%
\itemstyle{R\'ef\'erences~:}
N�ant.
